\section{Biased projections and simultaneous normalization}\label{sec:02-projections}

We now construct a family of projections whose means follow a prescribed curve in the cone. Its derivative field will provide both the Jacobians in the entropy argument and the matrix fields in the quadratic estimates. The final step of this section chooses coordinates and a Gaussian covariance simultaneously.

\subsection{The biased projection field}
For now fix the cone data satisfying the standing assumptions of Section~\ref{sec:01-cones}, and a symmetric matrix \(T\) with \(t_-I\le T\le t_+I\). Set
\[
\Sigma=I+T,\qquad Z=\Sigma^{1/2}G,
\]
where \(G\) is standard Gaussian in \(\mathbb R^m\). Thus every eigenvalue of \(\Sigma\) lies in \([.978,1.064]\). Expectations refer to \(G\), unless another variable is specified. For matrix fields use
\[
\tau B=\frac1m\mathbb E\operatorname{tr}B,\qquad
\tau_\Lambda B=\tau(\Lambda B),\qquad
\langle B,E\rangle_\Lambda=\tau_\Lambda(B\circ E),\qquad
B\circ E=\tfrac12(BE+EB).
\]
Here \(B,E\) are symmetric and \(\Lambda\) is a deterministic symmetric weight; the same notation applies to deterministic matrices. Transposition and cyclicity give \(\operatorname{tr}(\Lambda BE)=\operatorname{tr}(\Lambda EB)\). We suppress the subscript \(I\). For arbitrary square matrix fields write \(\|B\|_2^2=\tau(B^{\mathsf T}B)\), and put \([B,E]=BE-EB\).

Let \(\Pi_C\) denote Euclidean nearest-point projection onto \(C\), and put \(a(z)=\phi(z)+zp(z)\), so that \(a>0\) and \(a'=p\).

\begin{lemma}[Biased projections]\label{lem:biased-projections}
For every \(z\in\mathbb R\) there is a unique vector \(\xi_z\) such that
\[
X_z=\Pi_C(Z+\xi_z),\qquad Y_z=\Pi_D(-Z-\xi_z),\qquad
\mathbb EX_z=a(z)U.\tag{5}\label{eq:5}
\]
The map \(z\mapsto\xi_z\) is smooth. The a.e. derivative
\(P_z=D\Pi_C(Z+\xi_z)\) is a symmetric positive contraction, and
\[
X_z-Y_z=Z+\xi_z,\qquad X_z\cdot Y_z=0,\qquad
(\mathbb EP_z)\xi_z'=p(z)U.
\]
Define
\[
B(z)=\frac1m U\cdot\mathbb EY_z,\quad r_1(z)=-B'(z),\quad
h(z)=\tau_\Sigma P_z,\quad s_0=\tau\Sigma.
\]
Then
\[
\begin{gathered}
h(z)=\frac1m\mathbb E(X_z\cdot Z)
=\frac1m\mathbb E|X_z-\mathbb EX_z|^2+a(z)B(z),\\
h'=pB+ar_1,\qquad r_1>0.
\end{gathered}\tag{6}\label{eq:6}
\]
\end{lemma}
\begin{proof}
Moreau's cone decomposition, with the positive-dual convention, gives
\(\Pi_C(w)-\Pi_D(-w)=w\) and orthogonality of the two terms
\cite{Moreau1962,Soltan2019}. Here is the sign convention directly.
Writing $x=\Pi_C(w)$, the projection variational inequality is
$\langle w-x,c-x\rangle\le0$ for $c\in C$. Taking $c=0,2x$ gives
$\langle w-x,x\rangle=0$; the same inequality then gives $y=x-w\in D$.
For $d\in D$, $\langle-w-y,d-y\rangle=-\langle x,d\rangle\le0$,
so $y=\Pi_D(-w)$. Applying the projection inequality to two inputs
also gives
$|\Pi_C(w)-\Pi_C(w')|^2\le\langle\Pi_C(w)-\Pi_C(w'),w-w'\rangle$,
hence the projection is one-Lipschitz. It is the gradient of the convex function
\[
\Phi(w)=\tfrac12|\Pi_C(w)|^2
=\sup_{x\in C}\bigl(\langle w,x\rangle-\tfrac12|x|^2\bigr).
\]
Rademacher's theorem gives its a.e. differentiability
\cite{Rademacher1919,NekvindaZajicek1988}. A Lipschitz function is
absolutely continuous on coordinate lines; one-dimensional integration
by parts and Fubini's theorem identify these derivatives with the weak
derivatives. The convex-gradient representation therefore makes the
a.e. derivative symmetric and positive semidefinite, and the Lipschitz
bound makes it at most \(I\). Symmetry can equivalently be seen from commutation of weak mixed derivatives.

For fixed \(a>0\), minimize
\(\Psi_a(\xi)=\mathbb E\Phi(Z+\xi)-aU\cdot\xi\).
Write the orthogonal cone decomposition of \(\xi\) as \(\xi=x-y\). Since \(U\) is interior to \(C\), there is \(c>0\) with \(U\cdot y\ge c|y|\) for \(y\in D\). Jensen's inequality gives
\[
\Psi_a(\xi)\ge\tfrac12|x|^2-a|U||x|+ac|y|.
\]
This lower bound is coercive. Gaussian convolution makes \(\Psi_a\) smooth, with Hessian \(\mathbb E D\Pi_C(Z+\xi)\). This Hessian is strictly between \(0\) and \(I\): the nondegenerate Gaussian assigns positive probability to open subsets of \(\operatorname{int}C\), where the derivative is \(I\), and of \(-\operatorname{int}D\), where it is zero. Thus the minimizer is unique, and its stationarity condition is \(\mathbb EX=aU\). The implicit function theorem proves smooth dependence on \(z\) and the displayed formula for \(\xi_z'\).

Choose a measurable bounded symmetric-contraction representative for the derivative where it is undefined. For each layer these exceptional inputs are Gaussian-null; Fubini's theorem makes the choice immaterial to all layer integrals. Such a representative can be chosen jointly measurably by using difference quotients for the projection.

Gaussian integration by parts gives \(\mathbb E(X_z\cdot Z)=\mathbb E\operatorname{tr}(\Sigma P_z)\). Since \(\mathbb EY_z=aU-\xi_z\), orthogonality gives
\[
\mathbb E(X_z\cdot Z)=\mathbb E|X_z|^2-aU\cdot\xi_z.
\]
Subtracting \(|\mathbb EX_z|^2\) proves the variance expression in Equation~\eqref{eq:6}. Moreover, \(\nabla|\Pi_C|^2=2\Pi_C\), so
\[
\frac d{dz}\mathbb E|X_z|^2=2aU\cdot\xi_z',\qquad
h'=\frac1m(aU\cdot\xi_z'-pU\cdot\xi_z)=pB+ar_1.
\]
Finally,
\[
r_1=\frac p m U\cdot\bigl((\mathbb EP_z)^{-1}-I\bigr)U>0.
\]
No commutation of \(\Sigma\) with \(\mathbb EP_z\) is used.
\end{proof}

\subsection{Tails and layer integrals}
We need estimates that allow integration over all layers, despite possible nonsmoothness of the cone boundary. In Lemma~\ref{lem:projection-tails}, constants are locally uniform over invertible linear transformations of the fixed cone data, and uniform over the spectral box for \(T\). They may depend on the dimension.

\begin{lemma}[Gaussian layer tails]\label{lem:projection-tails}
For every finite \(r\ge1\), there are \(c_r,C_r>0\) such that
\[
\begin{aligned}
\|P_z\|_{L^r}+\|X_z\|_{L^r}&\le C_re^{-c_rz^2}&& (z\le0),\\
\|I-P_z\|_{L^r}+\|Y_z\|_{L^r}&\le C_re^{-c_rz^2}&& (z\ge0).
\end{aligned}
\]
Here any fixed finite-dimensional matrix norm can be used. Also \(B(z)\le C(1+|z|)\) for \(z\le0\). The positive measure \(\mu=h'(z)\,dz\) has total mass \(s_0\) and all polynomial moments. The measures with densities \(ar_1\) and \(pr_1\) have all polynomial moments as well.
\end{lemma}
\begin{proof}
Interiority and Equation~\eqref{eq:5} give
\[
\mathbb E|X_z|\le Ca(z),\qquad \mathbb E|Y_z|\le CB(z).
\]
For any one-Lipschitz vector function \(F\),
\(|F(Z)|\le\mathbb E|F(Z)|+\mathbb E|Z|+|Z|\); apply this to both projections. For \(s\ge0\),
\[
a(-s)=\phi(s)\int_0^\infty t e^{-st-t^2/2}\,dt
\asymp\frac{\phi(s)}{(1+s)^2},\qquad
p(-s)\le C\frac{\phi(s)}{1+s}.
\]
The lower bound follows by restricting the integral to \([0,(1+s)^{-1}]\); for the upper bound omit one or the other exponential factor. On \([0,\infty)\), \(a(z)\asymp1+z\).

For negative \(z\), split \(\mathbb E(X_z\cdot Z)\) at \(|Z|=C_0(1+|z|)\). The first part is bounded by \(C(1+|z|)a(z)\). In the second part the pointwise Lipschitz bound and Gaussian tails give the same bound when \(C_0\) is sufficiently large. Hence
\[
0\le h(z)\le C(1+|z|)a(z),\qquad B(z)\le h(z)/a(z)\le C(1+|z|).
\]
Since \(0\le P_z\le I\) and \(\Sigma\ge.978I\), this gives Gaussian decay of every finite moment of \(P_z\). The first moment of \(X_z\) is Gaussian-small, while the Lipschitz bound gives bounded moments of arbitrarily high order on the negative half-line. Interpolation gives the asserted decay of its \(L^r\) norms.

For positive \(z\), Equation~\eqref{eq:6} gives \(aB\le h\le s_0\), and thus
\[
|\mathbb EY_z|\le Cs_0/a(z),\qquad \xi_z=a(z)U-\mathbb EY_z.
\]
A fixed interior ball about \(U\) shows that \(\xi_z\) is a distance at least \(cz\) from the complement of \(C\) for all sufficiently large \(z\). On \(|Z|<cz\) the projection is the identity, so \(Y_z=0\) and \(P_z=I\) almost surely. Gaussian tails and the pointwise Lipschitz bound prove the remaining estimates.

The monotone function \(h\) tends to \(0\) and \(s_0\) at the two ends, with Gaussian tails. Integrating these tail functions proves the mass and moment assertions for \(\mu\). Finally \(ar_1\le h'\), and \(p/a\) has polynomial growth, proving the corresponding assertions for the other two densities.
\end{proof}

For a smooth scalar test \(f\) as in Section~\ref{sec:01-cones}, define
\[
\begin{gathered}
H_f=\int_{\mathbb R}(p(z)I-P_z)f'(z)\,dz,\qquad \iota(z)=z,\\
A=H_\iota,\qquad M_i=\mathbb E G_iA,\qquad
L=\sum_{i=1}^mG_iM_i,\qquad R=A-L.
\end{gathered}\tag{7}\label{eq:7}
\]
Lemma~\ref{lem:projection-tails} makes these layer integrals convergent in every finite \(L^r\), and all the displayed matrix fields have all finite moments.

\begin{lemma}[The Gaussian Sobolev interface]\label{lem:projection-sobolev}
For every smooth \(f\) with polynomially bounded derivatives, the centered vector field
\[
W_f=\int_{\mathbb R}(p(z)Z-X_z+\mathbb EX_z)f'(z)\,dz
\]
converges in Gaussian \(W^{1,2}\), and its weak Jacobian in \(G\) is
\(D_GW_f=H_f\Sigma^{1/2}\).
\end{lemma}
\begin{proof}
At negative layers the integrand tends to zero in \(L^2\) at a Gaussian rate. At positive layers use
\[
X_z-\mathbb EX_z=Z+Y_z-\mathbb EY_z
\]
to write it as \((p(z)-1)Z-Y_z+\mathbb EY_z\), with the same decay. The weak Jacobian of the integrand is \((p(z)I-P_z)\Sigma^{1/2}\), whose \(L^2\) norm has integrable Gaussian tails as well. Thus the truncated integrals and their weak derivatives are Cauchy in \(L^2\); closedness of the weak derivative gives the claim. Each integrand is centered. This argument differentiates the Lipschitz projections, not their derivative fields \(P_z\).
\end{proof}

For comparison, when \(C\) is the nonnegative orthant, \(U\) is the all-ones vector, and \(T=0\), one has \(\xi_z=zU\) and
\[
P_z=\operatorname{diag}(1_{\{-G_i\le z\}}),\qquad
A=L=\operatorname{diag}(-G_i).
\]
Indeed \(\int_{\mathbb R}(p(z)-1_{\{v_0\le z\}})\,dz=v_0\). In the general case the family \(P_z\) is neither assumed nested nor assumed to consist of projections. The field \(L\) is its deterministic-coefficient, degree-one Gaussian comparison.

\subsection{Choosing coordinates and covariance}
The remaining task is to arrange both a mean balance for \(A\) and a matrix equation for \(T\).

\begin{proposition}[Simultaneous normalization]\label{prop:normalization}
A linear transformation of the cone data and a symmetric \(T\) in the prescribed spectral box can be chosen so that
\[
\begin{gathered}
\mathbb EA=0,\qquad \lambda T^2+\mathbf C\circ T+\mathbf K=0,\qquad
\mathbf H:=-\mathbf C-2\lambda T\ge m_*I,\\
\mathbf C=\mathbb E\mathsf C(L),\qquad \mathbf K=\mathbb E\mathsf K(L).
\end{gathered}\tag{8}\label{eq:8}
\]
Matrix functions here and below use spectral calculus for symmetric matrices.
\end{proposition}
\begin{proof}
We use the scalar bounds in Proposition~\ref{prop:scalar-input}. For a
symmetric matrix \(Q\), transform the original data by
\[
 C_Q=e^QC_{\rm in},\qquad D_Q=e^{-Q}D_{\rm in},\qquad
 U_Q=e^QU_{\rm in},\qquad V_Q=e^{-Q}V_{\rm in}.
\]
Lemma~\ref{lem:cone-reduction} preserves the Laplace product and
\(U_Q\cdot V_Q=m\). Let \(T\) vary over its spectral box. For each
pair \((Q,T)\), construct the biases and fields as above; we suppress
their dependence on \(Q,T\) in the notation. We first establish
continuity, then construct a continuous self-map and exclude its fixed
points on a sufficiently large \(Q\)-sphere.

\paragraph{Continuity.}
Under convergent invertible transformations the cone projections converge uniformly on compact sets. To see this, projected points are bounded by the norm of the input; subsequential limits belong to the limiting cone and minimize the limiting distance, so uniqueness identifies the limit. Equicontinuity upgrades the conclusion to compact-uniform convergence. The coercive lower bound in Lemma~\ref{lem:biased-projections} is locally uniform, so the biases also converge at every fixed layer.

There is no need to assert pointwise convergence of the derivatives of these projections. Gaussian integration by parts instead gives
\[
\begin{aligned}
(\mathbb EP_z)\Sigma^{1/2}&=\mathbb E X_zG^{\mathsf T},\\
(\mathbb E G_iP_z)\Sigma^{1/2}
&=\mathbb E X_z(G_iG^{\mathsf T}-e_i^{\mathsf T}),
\end{aligned}
\]
where \(e_i\) is the \(i\)-th coordinate vector. The pointwise Lipschitz bounds dominate these expectations locally in the parameters. Lemma~\ref{lem:projection-tails}, also with a Gaussian factor \(G_i\) by H\"older's inequality, permits integration over layers. Hence \(\mathbb EA\) and all \(M_i\) are continuous. Their local boundedness and the polynomial growth of the profiles then give continuity of \(\mathbf C,\mathbf K\).

\paragraph{The covariance update.}
Put \(\Delta=\mathbf C^2-4\lambda\mathbf K\), and set
\[
T_{\rm new}=\frac{-\mathbf C-\sqrt\Delta}{2\lambda}.
\]
The scalar inequalities, spectral calculus, and expectation imply
\[
\begin{gathered}
\Delta\ge m_*^2I,\\
(\mathbf C+2\lambda t_+I)^2\le\Delta
\le(\mathbf C+2\lambda t_-I)^2.
\end{gathered}
\]
For the first inequality use \((\mathbf C+.37I)^2\ge0\). For the other two, expand the squares and use the corresponding endpoint inequalities in Proposition~\ref{prop:scalar-input}. Moreover \(\mathbf C+2\lambda t_\pm I<0\), since \(\mathbf C\le-.341I\) and \(2\lambda t_+=.0832\).

The L\"owner--Heinz inequality states that positive square root is order
preserving, including for noncommuting matrices
\cite[Theorem~2.3]{Chansangiam2013}. Indeed its integral representation uses the order-preserving matrices \(B(B+tI)^{-1}\). Taking square roots in the preceding sandwich therefore gives
\[
-\mathbf C-2\lambda t_+I\le\sqrt\Delta
\le-\mathbf C-2\lambda t_-I,
\]
so \(t_-I\le T_{\rm new}\le t_+I\). The positive discriminant gap also gives continuity of this update.

\paragraph{A uniform inward estimate.}
We claim that
\(\operatorname{tr}(Q\mathbb EA)<0\) whenever \(\sigma=\|Q\|_{\rm Frob}\) is sufficiently large, uniformly over its eigenbasis and over \(T\) in the box. All constants in this paragraph depend only on the original cone data and the fixed dimension, not on \(Q,T\).

Diagonalize \(Q\), write its eigenvalues as \(d_j\), and put \(h_j(z)=\mathbb E(P_z)_{jj}\). Applying the original interior-vector bounds to \(e^{-Q}X_z\) and \(e^QY_z\) gives
\[
\mathbb E|(X_z)_j|\le Ca(z)e^{d_j},\qquad
\mathbb E|(Y_z)_j|\le CB(z)e^{-d_j}.
\]
Gaussian integration by parts in this basis gives
\(h_j=\mathbb E[(X_z)_j(\Sigma^{-1}Z)_j]\), and the analogous identity for \(1-h_j\) uses \(-Y_z\). The covariance bounds and the one-Lipschitz coordinate bounds imply, for every cutoff \(J\ge1\),
\[
\begin{aligned}
h_j&\le CJ\mathbb E|(X_z)_j|+Ce^{-cJ^2},\\
1-h_j&\le CJ\mathbb E|(Y_z)_j|+Ce^{-cJ^2}.
\end{aligned}
\]
For example, bound \(|(X_z)_j|\) by its absolute mean plus \(\mathbb E|Z|+|Z|\), and split the Gaussian integral at \(|Z|=J\). A tail term containing the absolute mean is absorbed by the first term. These bounds require no commutation of \(Q\) and \(\Sigma\).

Set \(A_j=\int_{\mathbb R}(p-h_j)\,dz\) and
\(\varepsilon=1/\sqrt m\). The quantity to be made negative is
\(\operatorname{tr}(Q\mathbb EA)=\sum_jd_jA_j\).
Thus a negative eigenvalue requires a lower bound on \(A_j\), whereas
a positive eigenvalue requires an upper bound. We obtain a uniform bound
for each sign and then improve it when \(|d_j|\ge\varepsilon\sigma\).

If \(d_j\le0\), use \(J=1+|z|\) on the negative half-line in the
estimate for \(h_j\). The bound
\(\mathbb E|(X_z)_j|\le Ca(z)e^{d_j}\le Ca(z)\) makes its integral
bounded independently of \(Q,T\). On the positive half-line,
\(p-h_j\ge p-1\), whose integral is finite. Together these give
\(A_j\ge-C\).

If in addition \(d_j\le-\varepsilon\sigma\), the same estimate with
a sufficiently large fixed cutoff gives \(h_j\le1/4\) on
\([0,e^{\varepsilon\sigma/2}]\) for large \(\sigma\): here
\(a(z)e^{d_j}\le C(1+e^{\varepsilon\sigma/2})e^{-\varepsilon\sigma}\)
tends to zero. On this interval \(p-h_j\ge1/4\). The complementary
integrals retain their uniform lower bound, so
\[
A_j\ge\tfrac14e^{\varepsilon\sigma/2}-C.
\]

For positive eigenvalues we first note the uniform integral bound
\[
\int_0^\infty B\,dz\le2\int_0^\infty pB\,dz
\le2\int_0^\infty h'\,dz\le2s_0.
\]
On \([0,e^{2\sigma}]\), use \(J=C_1\sqrt\sigma\), with \(C_1\) large enough that the integrated Gaussian error is negligible. This bounds the integral of \(1-h_j\) there by \(C\sqrt\sigma e^{-d_j}+o(1)\).

The remaining tail must be controlled despite the changing cone. The transformed and scaled image of a fixed interior ball about the original \(U\) contains a ball about \(aU_Q\) of radius at least \(cae^{-\sigma}\). On the other hand,
\[
|\mathbb EY_z|\le Ce^\sigma B(z)\le Ce^\sigma/a(z).
\]
For \(z\ge e^{2\sigma}\), this error is negligible relative to the radius. Thus \(\xi_z=aU_Q-\mathbb EY_z\) has interior distance at least \(cze^{-\sigma}\), and
\[
1-h_j(z)\le C e^{-cz^2e^{-2\sigma}},\qquad
\int_{e^{2\sigma}}^\infty(1-h_j)\,dz
\le Ce^\sigma\int_{e^\sigma}^\infty e^{-cu^2}\,du=o(1).
\]
Combining these bounds with \(\int_{-\infty}^0p<\infty\), for every \(d_j\ge0\) we have
\[
A_j\le C+C\sqrt\sigma e^{-d_j}.
\]
If \(d_j\ge\varepsilon\sigma\), choose \(\rho>0\) with \(\rho^2/2<\varepsilon\). The lower bound on \(a\) gives
\[
\sup_{-\rho\sqrt\sigma\le z\le0}\frac{e^{-d_j}}{a(z)}
\le C(1+\sigma)e^{-(\varepsilon-\rho^2/2)\sigma}\longrightarrow0.
\]
For every real \(z\), the bound \(aB\le h\le s_0\) gives
\(B(z)\le s_0/a(z)\). A large fixed cutoff in the estimate for
\(1-h_j\) consequently gives
\(h_j\ge3/4\) throughout this interval. Its contribution to \(A_j\)
is at most \(-\rho\sqrt\sigma/4\), since \(p\le1/2\) there.
On the rest of the negative half-line use \(p-h_j\le p\); its
integral is bounded by \(\int_{-\infty}^0p\). On the positive
half-line use \(p-h_j\le1-h_j\) and the preceding bound
\(C\sqrt\sigma e^{-d_j}+o(1)\), which is bounded when
\(d_j\ge\varepsilon\sigma\). Hence
\(A_j\le-c\sqrt\sigma+C\).

At least one eigenvalue has \(|d_j|\ge\varepsilon\sigma\). A positive such eigenvalue contributes at most \(-c\sigma^{3/2}+O(\sigma)\) to \(\sum_jd_jA_j\); a negative such eigenvalue contributes a still larger negative amount in magnitude. Every remaining negative eigenvalue contributes at most \(C|d_j|\), and every remaining positive eigenvalue contributes at most
\(Cd_j+C\sqrt\sigma d_je^{-d_j}=O(\sigma)\).
There are only \(m\) terms. Hence \(\operatorname{tr}(Q\mathbb EA)=\sum_jd_jA_j<0\) for all sufficiently large \(\sigma\), as claimed.

\paragraph{The fixed point.}
Take a closed Frobenius ball of such a radius and the closed spectral box for \(T\). On their product update \(T\) as above and update \(Q\) by nearest-point projection of \(Q+\mathbb EA\) back to the ball. This is a continuous self-map of a finite-dimensional compact convex set. Brouwer's fixed point theorem \cite{Brouwer1911} applies. A fixed point on the \(Q\)-sphere would require \(\mathbb EA\) to be a nonnegative multiple of \(Q\), contradicting the inward estimate. At an interior fixed point, \(\mathbb EA=0\).

Finally write \(S=\sqrt\Delta\). At the fixed point \(T=(-\mathbf C-S)/(2\lambda)\), and direct expansion gives
\[
\lambda T^2+\mathbf C\circ T+\mathbf K
=\frac{S^2-\mathbf C^2}{4\lambda}+\mathbf K=0.
\]
The mixed products cancel without any commutation assumption. Also \(\mathbf H=S\ge m_*I\), proving Equation~\eqref{eq:8}.
\end{proof}

We henceforth use the fixed choices of Proposition~\ref{prop:normalization}. In particular \(T\) and all \(M_i\) are deterministic as functions of the Gaussian input. The balance \(\mathbb EA=0\) makes \(R=A-L\) orthogonal to Gaussian degrees zero and one.
